% Created 2024-01-05 vie 19:11
% Intended LaTeX compiler: pdflatex
\documentclass[11pt]{article}
\usepackage[utf8]{inputenc}
\usepackage[T1]{fontenc}
\usepackage{graphicx}
\usepackage{longtable}
\usepackage{wrapfig}
\usepackage{rotating}
\usepackage[normalem]{ulem}
\usepackage{amsmath}
\usepackage{amssymb}
\usepackage{capt-of}
\usepackage{hyperref}
\usepackage{minted}

\usepackage{parskip}
\usepackage[utf8]{inputenc} %% For unicode chars
\usepackage{placeins}
\usepackage[margin=2.50cm]{geometry}
\usepackage[T1]{fontenc}
\usepackage{mathpazo}
\linespread{1.05}
\usepackage[scaled]{helvet}
\usepackage{courier}
\hypersetup{colorlinks=true,linkcolor=blue}
\RequirePackage{fancyvrb}
\DefineVerbatimEnvironment{verbatim}{Verbatim}{fontsize=\small,formatcom = {\color[rgb]{0.5,0,0}}}
\usepackage{mdframed}
\BeforeBeginEnvironment{minted}{\begin{mdframed}}
\AfterEndEnvironment{minted}{\end{mdframed}}
\setcounter{secnumdepth}{3}
\date{}
\title{}
\hypersetup{
 pdfauthor={},
 pdftitle={},
 pdfkeywords={},
 pdfsubject={},
 pdfcreator={Emacs 28.1 (Org mode 9.6.12)}, 
 pdflang={Spanish}}
\begin{document}

\setcounter{tocdepth}{3}
\tableofcontents

\setlength\parindent{10pt}

\section{PHP Functions}
\label{sec:org19b1d5a}
Functions are one of the most essential constructs in programming. They
encapsulate code for specific tasks, enabling code reuse, logical
organization, and the simplification of complex operations. In this
chapter, we'll delve into defining and calling functions in PHP,
managing data passed to and from these functions, and using some of the
many built-in PHP functions.

\subsection{Defining and Calling Functions}
\label{sec:org2eef0dc}
A function is a designated block of code designed to perform a
particular task. In PHP, a function is defined using the \texttt{function}
keyword, followed by a unique function name and a set of parentheses.
Optional parameters can be specified within these parentheses, and these
act as input variables for the function. The function's body, enclosed
within curly braces, contains the code executed upon the function's
invocation.

Here's an example of a simple function definition that outputs a
greeting message:

\begin{minted}[]{php}
<?=
  function sayHello() {
      echo "Hello, world!";
  }
\end{minted}

To call this function, we use its name followed by parentheses. If the
function has parameters, the corresponding arguments are passed within
these parentheses. An argument is a value supplied to the function upon
its call.

Here's how we call the \texttt{sayHello} function:

\begin{minted}[]{php}
<?=
  sayHello(); // Outputs: Hello, world!
\end{minted}

\subsection{Function Arguments, Value and Reference, and Return Values}
\label{sec:orgc8734aa}
\subsubsection{Passing Arguments}
\label{sec:org3af60d5}
Parameters allow functions to accept inputs. These inputs are termed as
arguments when a function is called. Each argument corresponds to a
function's parameter by order.

Consider a function that greets a user by name:

\begin{minted}[]{php}
<?=
  function greet($name) {
      echo "Hello, $name!";
  }

  greet("Alice"); // Outputs: Hello, Alice!
\end{minted}

In this example, \texttt{"Alice"} is the argument passed to the \texttt{greet}
function.

\subsubsection{Value vs. Reference}
\label{sec:org5b35754}
In PHP, arguments can be passed by value or by reference. Passing by
value means the function works with a copy of the argument, whereas
passing by reference means the function works directly with the original
variable.

Example of passing by value:

\begin{minted}[]{php}
<?=
  function addOne($number) {
      $number++;
      echo $number;
  }

  $originalNumber = 5;
  addOne($originalNumber); // Outputs: 6
  echo $originalNumber; // Still outputs: 5
\end{minted}

In contrast, passing by reference:

\begin{minted}[]{php}
<?=
  function addOneByRef(&$number) {
      $number++;
  }

  $originalNumber = 5;
  addOneByRef($originalNumber);
  echo $originalNumber; // Outputs: 6
\end{minted}

The \texttt{\&} symbol in the parameter list indicates that the argument is
passed by reference.

\subsubsection{Return Values}
\label{sec:org0bfbf26}
Functions can return values using the \texttt{return} statement. This ends the
function's execution and optionally passes back a value.

For example:

\begin{minted}[]{php}
<?=
  function sum($a, $b) {
      return $a + $b;
  }

  $total = sum(5, 3);
  echo $total; // Outputs: 8
\end{minted}

\subsection{Built-in PHP Functions}
\label{sec:org611d081}
PHP offers an extensive collection of built-in functions that perform a
wide range of tasks.

\subsubsection{String Functions}
\label{sec:orgf5cfbeb}
String manipulation is a common requirement, and PHP provides many
functions to handle strings. Functions like \texttt{strlen}, \texttt{str\_replace}, and
\texttt{substr} are frequently used.

Example of string functions:

\begin{minted}[]{php}
<?=
  $text = "Hello, world!";
  echo strlen($text); // Outputs the length of the string
  echo str_replace("world", "PHP", $text); // Replaces 'world' with 'PHP'
\end{minted}

\subsubsection{Array Functions}
\label{sec:org920feff}
Arrays are essential in PHP, and functions like \texttt{array\_push}, \texttt{sort},
and \texttt{count} are integral for array manipulation.

Example of array functions:

\begin{minted}[]{php}
<?=
  $fruits = ["apple", "banana"];
  array_push($fruits, "orange"); // Adds 'orange' to the array
  sort($fruits); // Sorts the array
  echo count($fruits); // Outputs the number of elements in the array
\end{minted}

\subsubsection{Date and Time Functions}
\label{sec:org3761a0d}
Handling dates and times is another common requirement. Functions such
as \texttt{date} and \texttt{strtotime} are useful for formatting and manipulating
date and time.

Example of date and time functions:

\begin{minted}[]{php}
<?=
echo date("Y-m-d"); // Outputs the current date echo
strtotime("next Monday"); // Converts the time string to a Unix
timestamp
\end{minted}

\subsection{Summary}
\label{sec:orgf7414ed}
Functions in PHP are powerful tools for organizing and reusing code.
Understanding how to define, call, and manipulate data with functions is
key to effective PHP programming. Additionally, leveraging PHP's
built-in functions can greatly simplify many common tasks in web
development.

In the next sections, we will explore further into more advanced aspects
of PHP functions, including anonymous functions, variable scope, and
recursive functions. Stay tuned to deepen your understanding and enhance
your PHP programming skills.

\begin{itemize}
\item Anex I Below is a list of commonly used PHP built-in functions,
categorized by their functionality. This list is by no means
exhaustive, as PHP offers a vast library of built-in functions, but it
covers some of the most frequently used ones in various aspects of PHP
programming:
\end{itemize}

\section{Anex I: List of predefined and standard functions}
\label{sec:org72f8e3a}

\subsection{String Functions}
\label{sec:org5c623a9}
\begin{enumerate}
\item \texttt{strlen(\$string)} - Returns the length of a string.
\item \texttt{str\_replace(\$search, \$replace, \$subject)} - Replaces all occurrences
of a search string with a replacement string.
\item \texttt{strpos(\$haystack, \$needle)} - Finds the position of the first
occurrence of a substring in a string.
\item \texttt{strtolower(\$string)} - Converts a string to lowercase.
\item \texttt{strtoupper(\$string)} - Converts a string to uppercase.
\item \texttt{substr(\$string, \$start, \$length)} - Returns a part of a string.
\item \texttt{trim(\$string)} - Strips whitespace from the beginning and end of a
string.
\end{enumerate}

\subsection{Array Functions}
\label{sec:org2fdcfbe}
\begin{enumerate}
\item \texttt{array\_push(\&\$array, \$value1, \$value2, ...)} - Pushes one or more
elements onto the end of an array.
\item \texttt{sort(\&\$array)} - Sorts an array.
\item \texttt{count(\$array)} - Counts all elements in an array, or something in an
object.
\item \texttt{array\_merge(\$array1, \$array2, ...)} - Merges one or more arrays.
\item \texttt{array\_slice(\$array, \$offset, \$length)} - Extracts a slice of the
array.
\item \texttt{in\_array(\$needle, \$haystack)} - Checks if a value exists in an
array.
\item \texttt{array\_keys(\$array)} - Returns all the keys or a subset of the keys
of an array.
\end{enumerate}

\subsection{Mathematical Functions}
\label{sec:org756a3e7}
\begin{enumerate}
\item \texttt{abs(\$number)} - Returns the absolute value of a number.
\item \texttt{ceil(\$value)} - Rounds a number up to the nearest integer.
\item \texttt{floor(\$value)} - Rounds a number down to the nearest integer.
\item \texttt{rand(\$min, \$max)} - Generates a random integer.
\item \texttt{max(\$value1, \$value2, ...)} - Finds the highest value.
\item \texttt{min(\$value1, \$value2, ...)} - Finds the lowest value.
\item \texttt{round(\$val, \$precision)} - Rounds a float.
\end{enumerate}

\subsection{Date and Time Functions}
\label{sec:org88dda10}
\begin{enumerate}
\item \texttt{date(\$format, \$timestamp)} - Formats a local time/date.
\item \texttt{strtotime(\$time)} - Parses an English textual datetime into a Unix
timestamp.
\item \texttt{time()} - Returns the current time as a Unix timestamp.
\item \texttt{mktime(\$hour, \$minute, \$second, \$month, \$day, \$year)} - Gets Unix
timestamp for a date.
\item \texttt{date\_default\_timezone\_set(\$timezone\_identifier)} - Sets the default
timezone used by all date/time functions.
\end{enumerate}

\subsection{File Handling Functions}
\label{sec:org5c15f3d}
\begin{enumerate}
\item \texttt{fopen(\$filename, \$mode)} - Opens file or URL.
\item \texttt{fclose(\$handle)} - Closes an open file pointer.
\item \texttt{fwrite(\$handle, \$string)} - Writes to a file.
\item \texttt{file\_get\_contents(\$filename)} - Reads entire file into a string.
\item \texttt{file\_put\_contents(\$filename, \$data)} - Write data to a file.
\item \texttt{file\_exists(\$filename)} - Checks whether a file or directory exists.
\end{enumerate}

\subsection{Miscellaneous Functions}
\label{sec:org9a15a29}
\begin{enumerate}
\item \texttt{isset(\$var)} - Determine if a variable is set and is not NULL.
\item \texttt{empty(\$var)} - Determine whether a variable is empty.
\item \texttt{die(\$message)} - Equivalent to \texttt{exit()}, outputs a message and
terminates the current script.
\item \texttt{var\_dump(\$variable)} - Dumps information about a variable.
\item \texttt{gettype(\$var)} - Gets the type of a variable.
\end{enumerate}

This list serves as a starting point for exploring PHP's built-in
functions. Each function has its own set of parameters and behaviors,
which can be further explored in the PHP official documentation for more
detailed information and usage examples.
\end{document}
